\chapter{Results and Analysis}

\section{Data Summary}

This chapter reports the empirical results from the Bayesian logistic regression analysis. The empirical analysis uses UK accident data with a binary accident severity response, where slight accidents are modelled against the other severity categories. The model uses Bayesian logistic regression with the predictors and preprocessing steps described in Chapter 3.

Table \ref{tab:model_settings} summarises the computational settings used for the reported analysis. After preprocessing, the analysis contains 82,344 observations and 10 regression coefficients. The full-data samplers use 100,000 MCMC iterations, while the subset samplers for Consensus Monte Carlo use 50,000 MCMC iterations for each subset. The burn-in fraction is 0.2, and the Consensus Monte Carlo implementation uses $K=4$ subsets.

\begin{table}[htbp]
\centering
\caption{Model settings for the empirical analysis}
\label{tab:model_settings}
\begin{tabular}{lr}
\hline
Setting & Value \\
\hline
Observations after preprocessing & 82344 \\
Number of coefficients & 10 \\
Full-data MCMC iterations & 100000 \\
Burn-in fraction & 0.2 \\
Number of subsets, $K$ & 4 \\
Subset MCMC iterations & 50000 \\
\hline
\end{tabular}
\end{table}

The computational output does not contain a descriptive statistics table for the original variables or response distribution. Therefore, descriptive summaries by accident severity, road type, and speed limit are left for later inclusion.

TODO: add descriptive statistics table for the response and predictors when the output is available.

\section{Full-Data Metropolis--Hastings Results}

The full-data analysis was performed using Random Walk Metropolis--Hastings and Independent Metropolis--Hastings. The posterior summaries available for both full-data samplers include posterior mean, posterior median, posterior standard deviation, effective sample size, Monte Carlo standard error, 95\% credible interval, and odds-ratio summaries.

The regression coefficients are interpreted relative to the reference categories used in the design matrix. The exact reference categories for road type and speed limit should be stated once the design matrix coding is documented. TODO: identify the exact reference category for road type. TODO: identify the exact reference category for speed limit.

For logistic regression, exponentiating a coefficient gives the corresponding odds ratio. An odds ratio greater than one indicates higher posterior estimated odds of a slight accident relative to the relevant reference category or baseline condition, whereas an odds ratio less than one indicates lower posterior estimated odds. The odds-ratio summaries are available in the computational output, but separate odds-ratio tables are not added here to keep the main results tables compact.

\subsection{Random Walk Metropolis--Hastings}

The full-data Random Walk Metropolis--Hastings sampler has an acceptance rate of 0.39195. Table \ref{tab:full_rwmh_summary} reports selected posterior summaries for the regression coefficients. The posterior summaries are reported on the log-odds scale, while odds-ratio summaries are available in the computational output.

\begin{table}[htbp]
\centering
\caption{Full-data Random Walk Metropolis--Hastings posterior summaries}
\label{tab:full_rwmh_summary}
\scriptsize
\begin{tabular}{lrrrrr}
\hline
Parameter & Mean & Median & SD & Lower 95\% CrI & Upper 95\% CrI \\
\hline
Intercept & 1.9517 & 1.9507 & 0.0812 & 1.7955 & 2.1140 \\
Road type: one-way street & -0.2126 & -0.2129 & 0.0817 & -0.3719 & -0.0515 \\
Road type: roundabout & 0.1855 & 0.1856 & 0.0493 & 0.0900 & 0.2816 \\
Road type: single carriageway & -0.2267 & -0.2263 & 0.0390 & -0.3038 & -0.1517 \\
Road type: slip road & 0.5335 & 0.5306 & 0.1195 & 0.3062 & 0.7768 \\
Speed limit 2 & 0.1730 & 0.1742 & 0.0717 & 0.0296 & 0.3082 \\
Speed limit 3 & -0.0089 & -0.0088 & 0.0794 & -0.1645 & 0.1434 \\
Speed limit 4 & -0.2127 & -0.2124 & 0.0908 & -0.3922 & -0.0370 \\
Speed limit 5 & -0.3658 & -0.3648 & 0.0765 & -0.5195 & -0.2191 \\
Speed limit 6 & -0.2205 & -0.2193 & 0.0985 & -0.4144 & -0.0310 \\
\hline
\end{tabular}
\end{table}

For this sampler, the posterior means and posterior medians are close for all reported coefficients. The 95\% credible interval for speed limit 3 includes zero, while the remaining reported intervals do not include zero.

\subsection{Independent Metropolis--Hastings}

The full-data Independent Metropolis--Hastings sampler has an acceptance rate of 0.21504. Table \ref{tab:full_imh_summary} reports selected posterior summaries. The posterior means from the two full-data samplers are very close, suggesting that the two full-data Metropolis--Hastings implementations produce similar posterior summaries for this output.

\begin{table}[htbp]
\centering
\caption{Full-data Independent Metropolis--Hastings posterior summaries}
\label{tab:full_imh_summary}
\scriptsize
\begin{tabular}{lrrrrr}
\hline
Parameter & Mean & Median & SD & Lower 95\% CrI & Upper 95\% CrI \\
\hline
Intercept & 1.9500 & 1.9493 & 0.0797 & 1.7941 & 2.1116 \\
Road type: one-way street & -0.2130 & -0.2142 & 0.0813 & -0.3701 & -0.0539 \\
Road type: roundabout & 0.1867 & 0.1859 & 0.0486 & 0.0903 & 0.2824 \\
Road type: single carriageway & -0.2246 & -0.2248 & 0.0388 & -0.3014 & -0.1509 \\
Road type: slip road & 0.5370 & 0.5367 & 0.1187 & 0.3090 & 0.7746 \\
Speed limit 2 & 0.1729 & 0.1751 & 0.0713 & 0.0297 & 0.3101 \\
Speed limit 3 & -0.0091 & -0.0074 & 0.0789 & -0.1672 & 0.1425 \\
Speed limit 4 & -0.2108 & -0.2102 & 0.0904 & -0.3887 & -0.0360 \\
Speed limit 5 & -0.3652 & -0.3640 & 0.0769 & -0.5177 & -0.2172 \\
Speed limit 6 & -0.2189 & -0.2182 & 0.0987 & -0.4152 & -0.0245 \\
\hline
\end{tabular}
\end{table}

\section{Consensus Monte Carlo Results}

Consensus Monte Carlo was implemented using $K=4$ subsets. The computational output contains posterior summaries for Consensus Monte Carlo using Random Walk Metropolis--Hastings and Independent Metropolis--Hastings. The subset acceptance rates are reported in Table \ref{tab:acceptance_rates}.

\begin{table}[htbp]
\centering
\caption{Acceptance rates from the computational output}
\label{tab:acceptance_rates}
\begin{tabular}{lr}
\hline
Method & Acceptance rate \\
\hline
Full RWMH & 0.39195 \\
Full IMH & 0.21504 \\
CMC-RWMH subset 1 & 0.40850 \\
CMC-RWMH subset 2 & 0.39748 \\
CMC-RWMH subset 3 & 0.39106 \\
CMC-RWMH subset 4 & 0.38446 \\
CMC-IMH subset 1 & 0.02318 \\
CMC-IMH subset 2 & 0.14236 \\
CMC-IMH subset 3 & 0.03542 \\
CMC-IMH subset 4 & 0.01420 \\
\hline
\end{tabular}
\end{table}

\subsection{Consensus Monte Carlo using Random Walk Metropolis--Hastings}

Table \ref{tab:cmc_rwmh_summary} reports the posterior summaries for Consensus Monte Carlo using Random Walk Metropolis--Hastings. The CMC-RWMH posterior summaries are generally close to the full-data posterior summaries, although small differences are present for some coefficients.

\begin{table}[htbp]
\centering
\caption{Consensus Monte Carlo using Random Walk Metropolis--Hastings posterior summaries}
\label{tab:cmc_rwmh_summary}
\scriptsize
\begin{tabular}{lrrrrr}
\hline
Parameter & Mean & Median & SD & Lower 95\% CrI & Upper 95\% CrI \\
\hline
Intercept & 1.9485 & 1.9473 & 0.0807 & 1.7914 & 2.1085 \\
Road type: one-way street & -0.2128 & -0.2132 & 0.0824 & -0.3742 & -0.0522 \\
Road type: roundabout & 0.1790 & 0.1786 & 0.0497 & 0.0812 & 0.2757 \\
Road type: single carriageway & -0.2334 & -0.2334 & 0.0389 & -0.3092 & -0.1569 \\
Road type: slip road & 0.5123 & 0.5111 & 0.1219 & 0.2786 & 0.7513 \\
Speed limit 2 & 0.1830 & 0.1832 & 0.0726 & 0.0409 & 0.3250 \\
Speed limit 3 & 0.0010 & 0.0008 & 0.0799 & -0.1548 & 0.1576 \\
Speed limit 4 & -0.2131 & -0.2130 & 0.0918 & -0.3953 & -0.0300 \\
Speed limit 5 & -0.3590 & -0.3589 & 0.0780 & -0.5130 & -0.2076 \\
Speed limit 6 & -0.2222 & -0.2241 & 0.0987 & -0.4152 & -0.0275 \\
\hline
\end{tabular}
\end{table}

The CMC-RWMH subset acceptance rates range from 0.38446 to 0.40850. These values are close to the full-data RWMH acceptance rate of 0.39195.

\subsection{Consensus Monte Carlo using Independent Metropolis--Hastings}

The computational output contains posterior summaries for Consensus Monte Carlo using Independent Metropolis--Hastings. Therefore, the CMC-IMH results are reported in Table \ref{tab:cmc_imh_summary}. The subset acceptance rates for CMC-IMH are lower and more variable than those for CMC-RWMH, with values ranging from 0.01420 to 0.14236. These low subset acceptance rates suggest that the CMC-IMH sampler may require further tuning before strong conclusions are drawn from this method.

\begin{table}[htbp]
\centering
\caption{Consensus Monte Carlo using Independent Metropolis--Hastings posterior summaries}
\label{tab:cmc_imh_summary}
\scriptsize
\begin{tabular}{lrrrrr}
\hline
Parameter & Mean & Median & SD & Lower 95\% CrI & Upper 95\% CrI \\
\hline
Intercept & 1.9546 & 1.9575 & 0.0782 & 1.7897 & 2.1005 \\
Road type: one-way street & -0.1850 & -0.1838 & 0.0801 & -0.3413 & -0.0300 \\
Road type: roundabout & 0.1833 & 0.1832 & 0.0486 & 0.0871 & 0.2786 \\
Road type: single carriageway & -0.2300 & -0.2305 & 0.0387 & -0.3042 & -0.1542 \\
Road type: slip road & 0.5223 & 0.5200 & 0.1145 & 0.3000 & 0.7563 \\
Speed limit 2 & 0.1696 & 0.1687 & 0.0687 & 0.0376 & 0.3070 \\
Speed limit 3 & -0.0043 & -0.0062 & 0.0792 & -0.1587 & 0.1590 \\
Speed limit 4 & -0.2159 & -0.2191 & 0.0900 & -0.3886 & -0.0303 \\
Speed limit 5 & -0.3534 & -0.3553 & 0.0757 & -0.4969 & -0.2040 \\
Speed limit 6 & -0.2260 & -0.2265 & 0.0919 & -0.4070 & -0.0422 \\
\hline
\end{tabular}
\end{table}

\section{Comparison of Posterior Estimates}

Table \ref{tab:posterior_mean_comparison} compares posterior means across the four methods. The posterior means from the two full-data samplers are very close. The CMC-RWMH posterior means are also generally close to the full-data posterior means. The CMC-IMH posterior means are broadly similar for most parameters, although the result for the one-way street coefficient differs more visibly from the full-data estimates than the corresponding CMC-RWMH estimate.

\begin{table}[htbp]
\centering
\caption{Comparison of posterior means across methods}
\label{tab:posterior_mean_comparison}
\scriptsize
\begin{tabular}{lrrrr}
\hline
Parameter & Full RWMH & Full IMH & CMC-RWMH & CMC-IMH \\
\hline
Intercept & 1.9517 & 1.9500 & 1.9485 & 1.9546 \\
Road type: one-way street & -0.2126 & -0.2130 & -0.2128 & -0.1850 \\
Road type: roundabout & 0.1855 & 0.1867 & 0.1790 & 0.1833 \\
Road type: single carriageway & -0.2267 & -0.2246 & -0.2334 & -0.2300 \\
Road type: slip road & 0.5335 & 0.5370 & 0.5123 & 0.5223 \\
Speed limit 2 & 0.1730 & 0.1729 & 0.1830 & 0.1696 \\
Speed limit 3 & -0.0089 & -0.0091 & 0.0010 & -0.0043 \\
Speed limit 4 & -0.2127 & -0.2108 & -0.2131 & -0.2159 \\
Speed limit 5 & -0.3658 & -0.3652 & -0.3590 & -0.3534 \\
Speed limit 6 & -0.2205 & -0.2189 & -0.2222 & -0.2260 \\
\hline
\end{tabular}
\end{table}

The credible intervals are broadly similar across methods for most parameters. For all four methods, the 95\% credible interval for speed limit 3 includes zero. For the remaining reported coefficients, the available 95\% credible intervals do not include zero. This statement is descriptive and should not be interpreted as an assessment of practical importance without further model checking and contextual interpretation.

The effective sample size and Monte Carlo standard error are available in the posterior summaries. The full-data IMH output has larger effective sample sizes than the full-data RWMH output for the reported analysis. The CMC posterior summaries have smaller effective sample sizes than the corresponding full-data summaries in the computational output. These ESS values are not interpreted in isolation; they should be assessed together with trace plots, posterior plots, and other diagnostics.

No conclusion about computational efficiency is made because runtime values are not available in the computational output.

\section{Diagnostic Plots}

The computational output contains numerical posterior summaries and acceptance rates. However, trace plot files and posterior histogram or density plot files are not available.

\begin{figure}[htbp]
\centering
\fbox{\parbox{0.85\textwidth}{\centering TODO: insert trace plots for the full-data and Consensus Monte Carlo chains after the figure files are available.}}
\caption{Placeholder for trace plots}
\label{fig:trace_plots_placeholder}
\end{figure}

\begin{figure}[htbp]
\centering
\fbox{\parbox{0.85\textwidth}{\centering TODO: insert posterior histograms or density plots after the figure files are available.}}
\caption{Placeholder for posterior histograms or density plots}
\label{fig:posterior_plots_placeholder}
\end{figure}

Additional convergence diagnostics such as autocorrelation plots or potential scale reduction factors are not available in the computational output.

TODO: compute or export additional convergence diagnostics if required.

\section{Summary of Findings}

Based on the computational output, the full-data Random Walk Metropolis--Hastings and full-data Independent Metropolis--Hastings samplers produce very similar posterior means, posterior standard deviations, and 95\% credible intervals. This agreement provides a useful baseline for comparing the Consensus Monte Carlo results.

For Consensus Monte Carlo using Random Walk Metropolis--Hastings, the posterior summaries are generally close to the full-data posterior summaries. The CMC-RWMH subset acceptance rates are also close to the full-data RWMH acceptance rate. For Consensus Monte Carlo using Independent Metropolis--Hastings, posterior summaries are available, but the subset acceptance rates are lower and more variable than those from CMC-RWMH.

The results support a cautious descriptive conclusion that the Consensus Monte Carlo posterior summaries, especially those based on Random Walk Metropolis--Hastings, are broadly similar to the full-data posterior summaries for the computational output. However, no conclusion about computational efficiency is made because runtime values are not available. Further graphical diagnostics should be added after the trace plots and posterior histograms or density plots are exported.
